\begin{frame}{배치 정규화: 스킵 연결의 ``공범''}
  ResNet 논문은 스킵 연결과 함께 \kb{배치 정규화(batch
  normalization)}(Ioffe \& Szegedy, 2015)를 결합해 깊은 CNN의
  학습 안정성을 크게 높였다. 합성곱 층 \textbf{사이에} 넣는 층으로,
  각 \textbf{미니배치의 활성값 분포}를 평균 0·분산 1로 표준화한
  뒤, 학습 가능한 $\gamma$(스케일)·$\beta$(오프셋)로 다시
  곱해/더해 ``표준화가 너무 강하지 않게'' 조절:
  \[ \hat{x} = \frac{x - \mu_B}{\sqrt{\sigma_B^2 + \epsilon}},
     \qquad y = \gamma \hat{x} + \beta \]
  \begin{itemize}
    \item $\mu_B, \sigma_B^2$ = 해당 \textbf{미니배치 전체}
      활성값의 평균·분산, $\epsilon$ = 0 나눗셈 방지
      (예: $10^{-5}$)
    \item Week 9.2의 \textbf{초기화 문제를 \textbf{층마다}}
      해결: 100층을 쌓아도 활성값 분산이 층마다 달라지는
      (폭발/소실) 것을 \textbf{학습 중에도} 0·1로 되돌려,
      \textbf{초기화 값에 덜 민감하게} 깊게 쌓을 수 있게
  \end{itemize}
  \vskip0.4em
  ``깊게 쌓기''에 필요한 두 조건 --- 그래디언트가 살아서
  전달되고(스킵), 활성값이 적정 스케일로 유지되고(BN) --- 를
  \textbf{각각 독립적으로} 해결.
  {\footnotesize (10.3에서 BN 통계 처리가 전이학습의
  실전 함정이 되는 것 --- $\mu_B, \sigma_B$ 가 \textbf{학습 시}
  미니배치 기준이라 추론 시에는 \textbf{러닝된} 통계로
  대체해야 함.)}
\end{frame}
